% Autor szablonu: dr Roman Czapla
% Kompilacja: pdfLaTeX (dwukrotnie, aby zaktualizować spis treści).
% Wymagane logo uczelni: images/logo.png (ścieżka względem pliku głównego).
\documentclass[11pt,a4paper,oneside]{article}
\usepackage[utf8]{inputenc}
\usepackage[T1]{fontenc}
\usepackage[polish,provide=*]{babel}
\usepackage{lmodern,microtype}
\usepackage[margin=25mm,headheight=14pt]{geometry}
\usepackage{xcolor,graphicx,tabularx,booktabs,array,enumitem,fancyhdr,titlesec}
\usepackage[hidelinks]{hyperref}
\definecolor{accent}{HTML}{344B59}
\definecolor{hintcolor}{HTML}{58636A}
\hypersetup{pdfauthor={},pdftitle={Projekt inżynierski -- Dokumentacja projektowa}}
\setlength{\parindent}{0pt}
\setlength{\parskip}{6pt}
\renewcommand{\arraystretch}{1.3}
\setlength{\emergencystretch}{2em}
\setlist[itemize]{leftmargin=*,itemsep=2pt,topsep=3pt}
\titleformat{\section}{\Large\bfseries\color{accent}}{\thesection.}{0.5em}{}
\titleformat{\subsection}{\large\bfseries\color{accent}}{\thesubsection.}{0.5em}{}
\pagestyle{fancy}
\fancyhf{}
\fancyhead[L]{\small Projekt inżynierski}
\fancyhead[R]{\small Dokumentacja projektowa}
\fancyfoot[C]{\small\thepage}
\renewcommand{\headrulewidth}{0.3pt}
\newcolumntype{Y}{>{\raggedright\arraybackslash}X}

% Podpowiedzi są widoczne podczas wypełniania szablonu.
% Przed oddaniem: uzupełnij pola, zastąp przykłady własną treścią
% i zmień poniższy przełącznik na \podpowiedzifalse.
\newif\ifpodpowiedzi
\podpowiedzitrue
%\podpowiedzifalse
\newcommand{\podpowiedz}[1]{\ifpodpowiedzi\par\begingroup\small\color{hintcolor}\textit{Wskazówka: #1}\par\endgroup\fi}
\newcommand{\pole}[1]{\textnormal{[#1]}}

% Dane dokumentu -- uzupełnij w jednym miejscu.
% Ścieżkę logo można zmienić poniżej.
\newcommand{\pliklogo}{images/logo.png}
\newcommand{\logouczelni}{%
  \includegraphics[width=0.8\textwidth,height=40mm,keepaspectratio]{\pliklogo}%
}
\newcommand{\tytulprojektu}{\pole{Tytuł projektu}}
\newcommand{\rokakademicki}{\pole{2026/2027}}
\newcommand{\semestrprojektu}{\pole{zimowy / letni}}
\newcommand{\datadokumentu}{\pole{dd.mm.rrrr}}
\newcommand{\wersjadokumentu}{\pole{1.0}}
% Zespół dwuosobowy tylko za zgodą dyrekcji: zmień na \trzeciaosobafalse.
\newif\iftrzeciaosoba
\trzeciaosobatrue
\newcommand{\studentA}{\pole{Imię i nazwisko osoby 1}}
\newcommand{\studentB}{\pole{Imię i nazwisko osoby 2}}
\newcommand{\studentC}{\pole{Imię i nazwisko osoby 3}}
\newcommand{\albumA}{\pole{numer albumu}}
\newcommand{\albumB}{\pole{numer albumu}}
\newcommand{\albumC}{\pole{numer albumu}}
\newcommand{\prowadzacyA}{\pole{Stopień lub tytuł, imię i nazwisko}}
\newcommand{\prowadzacyB}{\pole{Stopień lub tytuł, imię i nazwisko}}

\begin{document}
\begin{titlepage}
\centering
{\large Uniwersytet Komisji Edukacji Narodowej w Krakowie\par}
{\normalsize Instytut Bezpieczeństwa i Informatyki\par}
\vspace{8mm}
\logouczelni\par
\vspace{10mm}
{\Large\bfseries\color{accent} PROJEKT INŻYNIERSKI\par}
\vspace{3mm}
{\large Dokumentacja projektowa\par}
\vspace{15mm}
{\LARGE\bfseries\tytulprojektu\par}
\vspace{15mm}
{\small ZESPÓŁ PROJEKTOWY\par}
\vspace{3mm}
\begin{tabular}{ll}
\studentA & Nr albumu: \albumA\\
\studentB & Nr albumu: \albumB
\iftrzeciaosoba\\
\studentC & Nr albumu: \albumC\fi
\end{tabular}
\vspace{12mm}

{\small PROWADZĄCY ZAJĘCIA PROJEKTOWE\par}
\vspace{3mm}
\prowadzacyA\par
\prowadzacyB\par
\vfill
Rok akademicki: \rokakademicki\quad Semestr: \semestrprojektu\par
Data dokumentu: \datadokumentu\quad Wersja: \wersjadokumentu\par
\vspace{6mm}
Kraków\par

\end{titlepage}

\tableofcontents
\podpowiedz{Dokument opisuje finalne rozwiązanie i umożliwia zrozumienie jego budowy oraz dalszą pracę nad projektem. Zakres uzgodnij z prowadzącymi: usuń podsekcje nieadekwatne do projektu i dodaj potrzebne. UML, baza danych i opis kodu nie są obowiązkowe dla każdego rodzaju projektu. Opis przebiegu pracy i indywidualnego wkładu należy do sprawozdania końcowego. Uzupełnij pola i wyłącz podpowiedzi przed oddaniem.}

\section{Szczegółowa dokumentacja projektowa}
\subsection{Zakres i identyfikacja rozwiązania}
\podpowiedz{Opisz przeznaczenie projektu, najważniejsze wymagania oraz wersję, której dotyczy dokumentacja. Wskaż, gdzie znajdują się pliki źródłowe i zasoby: repozytorium dla projektu programistycznego lub inne miejsce uzgodnione z prowadzącymi. Podaj commit, tag albo oznaczenie wersji plików.}
\textbf{Przeznaczenie i zakres:} \pole{opis.}\par
\textbf{Wersja projektu:} \pole{oznaczenie i data.}\par
\textbf{Materiały źródłowe:} \pole{lokalizacja i organizacja.}

\newpage
\subsection{Architektura i organizacja projektu}
\podpowiedz{Przedstaw główne elementy rozwiązania, ich odpowiedzialności i powiązania. Dla aplikacji opisz komponenty i przepływ danych; dla animacji 3D -- organizację scen, obiektów, zasobów i etapów przetwarzania. Wyjaśnij najważniejsze decyzje i zależności. Dodaj schemat, jeśli pomaga zrozumieć rozwiązanie.}
\pole{Budowa rozwiązania, struktura plików i uzasadnienie decyzji.}

\subsection{Modele i diagramy projektowe}
\podpowiedz{Dobierz diagramy do projektu. Dla oprogramowania mogą to być diagramy UML: klas, przypadków użycia, sekwencji, stanów lub komponentów. Dla innych projektów wykorzystaj np. schemat sceny, storyboard lub diagram procesu. Każdy rysunek opatrz podpisem i omów w tekście, stosując odwołania przez \texttt{\textbackslash ref}. Nie twórz diagramów wyłącznie w celu wypełnienia szablonu.}
\pole{Modele i diagramy z objaśnieniami.}
% Przykład rysunku -- odkomentuj po dodaniu pliku:
% \begin{figure}[htbp]
%   \centering
%   \includegraphics[width=0.9\linewidth]{images/architektura.png}
%   \caption{Architektura rozwiązania.}
%   \label{fig:architektura}
% \end{figure}
% Odwołanie: Na rysunku \ref{fig:architektura} przedstawiono ...

\subsection{Dane, baza danych i zasoby}
\podpowiedz{Jeśli projekt korzysta z bazy danych, opisz schemat, tabele, klucze, relacje, ograniczenia i istotne procedury. Dla innych danych lub zasobów określ formaty, strukturę, pochodzenie oraz sposób przygotowania i wykorzystania. Dla animacji mogą to być modele, tekstury, materiały i pliki dźwiękowe. Wskaż zasoby własne i zewnętrzne oraz warunki ich wykorzystania.}
\pole{Struktura danych lub zasobów i sposób ich wykorzystania.}

\subsection{Implementacja i najważniejsze rozwiązania}
\podpowiedz{Dla projektu programistycznego opisz istotne moduły, klasy, funkcje, interfejsy i algorytmy. Podaj znaczenie parametrów, wyniki, ograniczenia i obsługę błędów tam, gdzie jest to potrzebne. Nie wymaga się przepisywania całego kodu ani opisu każdej zmiennej. Dokumentację generowaną z kodu można dołączyć lub wskazać jako osobny materiał.}
\podpowiedz{Dla projektu bez kodu opisz odpowiednie metody wykonania, np. modelowanie, rigging, animację, oświetlenie, rendering i montaż. Podaj istotne ustawienia i uzasadnij ich dobór. Usuń wskazówki nieodnoszące się do projektu.}
\pole{Opis zastosowanych metod i kluczowych elementów wykonania.}

\newpage
\subsection{Mechanizmy AI w rozwiązaniu}
\podpowiedz{Ta podsekcja dotyczy AI stanowiącej część projektu. Opisz model lub usługę, wersję, zadanie, dane wejściowe i wyjściowe, integrację, konfigurację i ograniczenia. Jeśli rozwiązanie obejmuje uczenie, opisz dane, podział zbiorów, parametry i ocenę wyników. Dla usług zewnętrznych wskaż wymagane połączenia i zależności, bez ujawniania kluczy dostępowych. Jeśli AI tylko wspomagała pracę zespołu, jej wykorzystanie opisz w sprawozdaniu, a tę podsekcję usuń.}
\pole{Budowa i działanie mechanizmów AI -- jeśli występują.}

\subsection{Środowisko i odtworzenie projektu}
\podpowiedz{Opisz środowisko potrzebne do dalszej edycji, budowania lub odtworzenia rezultatu: system operacyjny, sprzęt, programy, biblioteki i ich wersje. Dla animacji wskaż np. wersję programu, dodatki, silnik renderujący i ustawienia eksportu. Szczegółową instrukcję instalacji i konfiguracji można wskazać w osobnej dokumentacji instalacyjno-konfiguracyjnej, jeśli jest wymagana. Nie powielaj tych samych informacji; podaj odwołanie do odpowiedniej sekcji.}
\begin{tabularx}{\textwidth}{@{}Y l Y@{}}
\toprule
Narzędzie / zależność & Wersja & Zastosowanie\\
\midrule
\pole{nazwa} & \pole{wersja} & \pole{rola w projekcie}\\
\pole{nazwa} & \pole{wersja} & \pole{rola w projekcie}\\
\bottomrule
\end{tabularx}

\textbf{Odtworzenie rezultatu:}
\pole{kroki edycji, budowania, renderowania lub innego przetwarzania; ewentualne odwołania do dokumentacji instalacyjno-konfiguracyjnej.}

\subsection{Testy i weryfikacja}
\podpowiedz{Powiąż sprawdzenia z wymaganiami. Opisz warunki, dane lub scenariusze, oczekiwane rezultaty, uzyskane wyniki i wnioski. Dobierz metody do charakteru projektu: testy oprogramowania, sprawdzenia prototypu, kontrola jakości animacji lub ocena mechanizmu AI. Wskaż ograniczenia przeprowadzonych sprawdzeń.}
\begin{tabularx}{\textwidth}{@{}Y Y Y@{}}
\toprule
Wymaganie / scenariusz & Sposób sprawdzenia & Wynik i wniosek\\
\midrule
\pole{wymaganie} & \pole{metoda i warunki} & \pole{wynik}\\
\pole{wymaganie} & \pole{metoda i warunki} & \pole{wynik}\\
\bottomrule
\end{tabularx}

\subsection{Ograniczenia i dalszy rozwój}
\podpowiedz{Wskaż znane błędy, ograniczenia techniczne, niespełnione wymagania i zależności wpływające na dalszą pracę. Oddziel stan istniejący od propozycji rozwoju.}
\pole{Ograniczenia i możliwe rozszerzenia.}

\subsection{Źródła i materiały referencyjne}
\podpowiedz{Wymień dokumentację, publikacje, biblioteki, modele i zasoby przywołane w dokumencie. Dla źródeł internetowych podaj tytuł, adres i datę dostępu. Wskaż licencje lub inne warunki wykorzystania materiałów zewnętrznych. Nie pozostawiaj fikcyjnych źródeł.}
\pole{Wykaz źródeł i materiałów referencyjnych.}
\end{document}
